% !TEX program = lualatex
% !TeX encoding = UTF-8
\PassOptionsToPackage{unicode}{hyperref}
\PassOptionsToPackage{bookmarksnumbered,bookmarksopen}{hyperref}
\documentclass[11pt,a4paper]{article}

% ============================================================
% إعدادات اللغة العربية (LuaLaTeX + polyglossia)
% ============================================================
\usepackage{polyglossia}
\setmainlanguage{arabic}
\setotherlanguage{english}

% خطوط عربية
\newfontfamily\arabicfont{Amiri}[Script=Arabic, Language=Arabic]
\newfontfamily\arabicfontsf{Amiri}[Script=Arabic, Language=Arabic]
\newfontfamily\arabicfonttt{Amiri}[Script=Arabic, Language=Arabic]
\newfontfamily\englishfont{Times New Roman}
\setmonofont{Latin Modern Mono}

% إزالة تحذيرات hyperref
\makeatletter
\let\@ensure@LTR\relax
\makeatother

% ============================================================
% حزم رياضيات وتنسيق (مطابقة للأصل)
% ============================================================
\usepackage{amsmath,amssymb,amsthm,mathtools,bm}
\usepackage{geometry}
\usepackage{enumitem}
\usepackage{siunitx}
\usepackage{booktabs}
\usepackage{array}
\usepackage{slashed}
\usepackage{graphicx}
\usepackage{caption}
\usepackage{subcaption}
\usepackage{braket}
\usepackage{tensor}
\usepackage{makecell}
\usepackage[most]{tcolorbox}
\usepackage{pgfplots}
\usepackage{float}
\usepackage{tikz}
\usepackage{multicol}
\usepackage{lipsum}
\usepackage{microtype}
\usepackage{hyperref}
\usepackage{longtable}
\usepackage{placeins}
\usepackage{xcolor}

\pgfplotsset{compat=1.18}
\usetikzlibrary{arrows.meta,positioning,shapes.geometric,fit,calc,
	decorations.pathreplacing,patterns,quotes,angles,
	shadows,decorations.markings}
\usetikzlibrary{shapes.geometric}

\geometry{margin=1in}
\setlength{\emergencystretch}{3em}
\sloppy

% بيئات النظريات
\newtheorem{theorem}{مبرهنة}[section]
\newtheorem{lemma}{تمهيدية}[section]
\newtheorem{axiom}{مسلَمة}[section]
\newtheorem{remark}{ملاحظة}[section]
\newtheorem{proposition}{اقتراح}[section]
\newtheorem{definition}{تعريف}[section]
\newtheorem{assumption}{افتراض}[section]
\newtheorem{corollary}{نتيجة طبيعية}[section]

% بيانات PDF
\hypersetup{
	pdftitle={الملحق D: تنبؤات YUCT البيولوجية مع الأسس الفيزيائية المتكاملة},
	pdfauthor={أليكسي ف. ياكوشيف},
	pdfsubject={كفاءة التنسيق، سلم الكتلة، القياس البيولوجي، معدلات الطفرات، القياس الأيضي، التزامن العصبي},
	pdfkeywords={YUCT, كفاءة التنسيق, تكميم الكتلة, قياس بيولوجي, سلم شامل},
	breaklinks=true,
	pdfencoding=auto,
	bookmarksnumbered=true,
	bookmarksopen=true,
	bookmarksopenlevel=1
}

% تعريف أوامر مختصرة
\newcommand{\Keff}{K_{\mathrm{eff}}}
\newcommand{\kappac}{\kappa_c}
\newcommand{\eps}{\varepsilon}
\newcommand{\betaf}{\beta}
\newcommand{\df}{d_{\!f}}
\newcommand{\Sodd}{S_{\mathrm{odd}}}
\newcommand{\Seven}{S_{\mathrm{even}}}
\newcommand{\Mpl}{M_{\mathrm{Pl}}}
\newcommand{\REL}{\mathcal{K}}

\begin{document}
	
	\title{\Huge\textbf{الملحق D. تنبؤات YUCT البيولوجية}\\[0.3cm]
		\Large نسخة مستعادة بالكامل مع جداول بيانات شاملة،\\[0.2cm]
		تحقق تجريبي موسع، وروابط صريحة مع BCLM/BCLM-EC}
	\author{أليكسي ف. ياكوشيف \\ 
		أبحاث ياكوشيف \\ \url{https://yuct.org/} \\ \url{https://ypsdc.com/}}
	\date{مارس 2026 \\ الإصدار 3.2 (جداول مستعادة وتكامل كامل)}
	
	\maketitle
	
	\begin{abstract}
		\noindent
		يقدم هذا الملحق تنبؤات بيولوجية قابلة للاختبار مستمدة من نظرية التنسيق الموحّد لياكوشيف (YUCT)، وهي الآن متكاملة بالكامل مع لاغرانجيان التنسيق البيولوجي (الملحق BCLM) وساعات التنسيق BCLM (الملحق BCLM-EC). نوضّح أن تكميم الكتلة نصف الصحيح للفرميونات (الملحق~Y) وكمّ التدرج \(q = (3/2)^{1/3}\) يمتدان إلى الأنظمة الحية، واضعًا معقدات البروتين، الفيروسات، البكتيريا، وخلايا كاملة على نفس الشبكة الإحداثية للجسيمات الأولية. قانون الخطأ الشامل \(\varepsilon \propto K_{\mathrm{eff}}^{-0.67}\) يحكم معدلات الطفرات (68 نوعًا فقاريًا، \(R^2 = 0.91\))، القياس الأيضي (1016 مشاهدة نباتية، أس \(0.68 \pm 0.04\))، التزامن العصبي، تراكم الطفرات الجسدية، الساعات فوق الجينية، تقييد السعرات، استشعار النصاب، وحركية تطوي البروتين. الدلالة الإحصائية المجتمعة تتجاوز \(p < 10^{-18}\). جميع جداول البيانات الموثقة سابقًا مستعادة ومربوطة بالإطار الرياضي. سلم التنسيق الكبير الموحّد يعرض جميع الأجسام المستقرة المعروفة واقعة على خط واحد بميل \(\ln q\). الإشارات المرجعية الصريحة إلى BCLM (كفاءات الكودونات، الأحماض الأمينية، الإنزيمات) وإلى BCLM-EC (ساعات التنسيق المبنية على المثيلة) تبرهن على الاتساق الكامل للقطاع البيولوجي من YUCT.
	\end{abstract}
	
	\vspace{0.2cm}
	\noindent
	\textbf{كلمات مفتاحية:} YUCT، عمق التنسيق، سلم الكتلة، تكميم نصف صحيح، كتل الفرميونات، القياس البيولوجي، معدل الأيض، معدل الطفرات، تزامن عصبي، طفرات جسدية، ساعات فوق جينية، تقييد السعرات، استشعار النصاب، تطوي البروتين، سلم شامل، تنبؤات قابلة للتفنيد.
	
	\vspace{0.1cm}
	\noindent
	\textbf{اقتباس:} ياكوشيف أ. ف. الملحق D. تنبؤات YUCT البيولوجية مع الأسس الفيزيائية المتكاملة (جداول مستعادة). 2026. DOI: \url{https://doi.org/10.5281/zenodo.18444598} (هذا المجلد).
	
	\tableofcontents
	\newpage
	
	\section{مقدمة}
	
	تقترح نظرية التنسيق الموحّد لياكوشيف (YUCT) أن جميع البنى المستقرة في الطبيعة — من الكواركات إلى المجرات — منظمة بمبدأ تنسيق شامل. في قلبها تقع الحلقة الجبرية \(S_{\mathrm{odd}} = 6/5\)، \(S_{\mathrm{even}} = 4/5\)، معطية النسبة الأساسية \(3/2 = S_{\mathrm{odd}}/S_{\mathrm{even}}\) وأس الخطأ \(\beta = 2/3\). قانون الخطأ الشامل \(\varepsilon = \kappa_c \alpha K_{\mathrm{eff}}^{-\beta}\) مع \(\kappa_c = 1/3\) تم التحقق منه عبر أكثر من 50 رتبة أسية في الطاقة والطول (الملحق~L).
	
	أحد التنبؤات الرئيسية هو أن جميع الكتل مكمّاة بخطوات نصف صحيحة لكمّ التدرج \(q = (3/2)^{1/3}\). نفس الكمّ الذي يعيد إنتاج كتل جميع اللبتونات والكواركات المشحونة بدقة دون في المئة (الملحق~Y) يحدد أيضًا كتل معقدات البروتين، الريبوسومات، وخلايا كاملة. في هذه النسخة المستعادة بالكامل، نربط صراحةً القياس البيولوجي بالأسس الفيزيائية لـ YUCT، مشيرين إلى نظاميات الكتلة النووية (الملحق~Y)، التداعيات الكونية (الملحق~\(\Omega\))، ولاغرانجيان التنسيق البيولوجي المُدخل حديثًا (BCLM، الملحق BCLM) مع ساعة BCLM فوق الجينية (BCLM-EC، الملحق BCLM-EC). نقدم أيضًا تصوّرًا موحدًا — سلم التنسيق الكبير الموحّد — الذي يعرض جميع الأجسام المستقرة المعروفة، من الإلكترون إلى الشجرة البالغة، واقعة على خط صلب واحد بميل \(\ln q\).
	
	علاوة على ذلك، تتضمن الوثيقة الآن جداول بيانات موسعة كانت جزءًا من النسخ السابقة لكنها أُغفلت سهوًا في المسودة السابقة. هذه الجداول تعرض البيانات التجريبية الأصلية (معدلات الطفرات، القياس الأيضي، التزامن العصبي، إلخ) إلى جانب كمياتها المشتقة من YUCT. التوافق بين هذه البيانات والنظرية ممتاز ومُكمّم إحصائيًا.
	
	\subsection{مبادئ أساسية}
	
	\begin{enumerate}
		\item \textbf{تنسيق قائم على القاموس:} الجزيئات البيولوجية ترث بروتوكولات تنسيق مسبقة الترميز (قواميس) تأسست عبر التطور.
		\item \textbf{تضخيم الرنين:} الأنظمة الجزيئية تُظهر تعزيزًا رنّانًا (\(R > 1\)) عندما تقترب نسب كتلها من قوى \(3/2\).
		\item \textbf{كفاءة قابلة للقياس:} كفاءة التنسيق \(K_{\mathrm{eff}}\) تحدد معدلات الخطأ، الشدة الأيضية، وطول العمر.
		\item \textbf{سلم كتلة شامل:} كتلة أي كيان منسق مستقر مقيدة بدرجات متقطعة على السلم \(m = m_e q^{N_f}\)، حيث \(N_f\) عدد صحيح أو نصف صحيح.
	\end{enumerate}
	
	\subsection{منهج التحقق الاستعادي}
	
	نتحقق من التنبؤات مقابل بيانات متاحة للعموم:
	\begin{itemize}
		\item \textbf{معدلات الطفرات:} طفرات الخط الجنسي عبر 68 نوعًا فقاريًا (Bergeron et al. 2023) \cite{Bergeron2023a}.
		\item \textbf{القياس الأيضي:} قاعدة بيانات الجذور النباتية العالمية (1016 مشاهدة) \cite{Yuan2020}.
		\item \textbf{التزامن العصبي:} بيانات تدريب شبكات عصبية مستزرعة \cite{PLOS2025b}.
		\item \textbf{تطوي البروتين:} قاعدة بيانات K‑Pro (1529 سجلًا) \cite{Turina2023}.
		\item \textbf{الطفرات الجسدية:} تحالف PCAWG \cite{PCAWG2020}.
		\item \textbf{الساعات فوق الجينية:} Horvath (2013) \& Lu et al. (2023) \cite{Horvath2013, Lu2023}.
		\item \textbf{تقييد السعرات:} تجربة CALERIE \cite{CALERIE2018, CALERIE2024}.
		\item \textbf{الخلايا الاصطناعية:} JCVI‑syn3.0 \cite{Hutchison2016, Pelletier2023}.
		\item \textbf{استشعار النصاب:} Miller \& Bassler (2001) \cite{Miller2001}, Papenfort \& Bassler (2016) \cite{Papenfort2016}.
	\end{itemize}
	
	\section{سلم الكتلة الشامل وأساسه الفيزيائي}
	
	\subsection{اشتقاق كمّ التدرج}
	
	قانون الخطأ الشامل يتبع من البعد الكسوري لمسارات التنسيق (الملحق~Y). جمع المتتاليات الهندسية على طبقات التنسيق الفردية والزوجية يعطي
	\[
	S_{\mathrm{odd}} = \frac{\beta}{1-\beta^2} = \frac{6}{5}, \quad
	S_{\mathrm{even}} = \frac{\beta^2}{1-\beta^2} = \frac{4}{5},
	\]
	مع \(\beta = 2/3\). النسبة \(3/2 = S_{\mathrm{odd}}/S_{\mathrm{even}}\) هي درجة الكتلة الأساسية بين طبقات التنسيق المتعاقبة. لأن الأبعاد المكانية الثلاثة تجمع الأخطاء بشكل مضاعف، فإن الخطوة الأولية على سلم الكتلة اللوغاريتمي هي الجذر التكعيبي: \(q = (3/2)^{1/3}\). وبالتالي، فإن كتلة أي جسم مستقر مقيدة بالمجموعة المتقطعة
	\begin{equation}
		m = m_e \cdot q^{N_f}, \qquad N_f \in \tfrac12\mathbb{Z},
		\label{eq:quantization}
	\end{equation}
	حيث \(m_e\) كتلة الإلكترون. هذا التنبؤ خالٍ من المعاملات الحرة؛ المُدخلات الوحيدة هي \(\beta = 2/3\) وكتلة الإلكترون كنقطة مرجعية.
	
	\subsection{الامتداد إلى النوى والجزيئات البيولوجية الكبيرة}
	
	نفس السلم ينظم النوى الذرية (الملحق~Y) والجزيئات البيولوجية الكبيرة (الملحق BCLM-EC، القسم~4). عمق التنسيق الفعّال لنواة بعدد كتلي \(A\) هو \(L_{\mathrm{eff}} \approx 96 - \frac13 \log_{3/2} A + \delta(Z,A)\). كتل النوى الخفيفة (\(^4\)He، \(^{12}\)C، \(^{16}\)O) تتوافق بدقة مع قيم \(N_f\) نصف الصحيحة. علاوة على ذلك، الكتل المبلَّغة حديثًا لمعقدات البروتين — يوبيكويتين، ريبوسوم، بروتيازوم، نوكليوسوم، معقد المسام النووي — جميعها تقع ضمن \(\pm 0.25\) من العقد نصف الصحيحة (انظر جدول~\ref{tab:protein_masses} في الملحق BCLM-EC). الاحتمال المشترك لمثل هذه المحاذاة بالصدفة هو \(< 10^{-4}\). هذا يؤسس لعالمية السلم من فيزياء الجسيمات إلى البيولوجيا الجزيئية.
	
	\subsection{سلم التنسيق الكبير الموحّد}
	
	الشكل~\ref{fig:grand_ladder} يعرض سلم الكتلة الشامل، المتضمن الآن الخلية الاصطناعية Syn3.0 والكائنات الحية الكبرى المتنبأ بها.
	
	\begin{figure}[H]
		\centering
		\begin{tikzpicture}
			\begin{axis}[
				width=0.95\textwidth, height=0.55\textwidth,
				xlabel={الدليل نصف الصحيح \(N_f\)},
				ylabel={\(\ln(m/m_e)\)},
				grid=major,
				legend pos=north west,
				legend cell align=left,
				legend style={font=\footnotesize},
				title={سلم التنسيق الموحّد: من الفرميونات إلى النظم البيئية},
				title style={font=\bfseries},
				xtick={0,50,...,400},
				minor xtick={0,10,...,400},
				ymin=-2, ymax=55,
				xmin=-10, xmax=410,
				]
				
				\addplot[domain=-20:420, samples=2, dashed, ultra thick, blue!60] 
				{x * 0.135155};
				\addlegendentry{النظرية: \(\ln m = N_f \ln q\) (لا معاملات حرة)}
				
				\addplot[only marks, mark=*, red, mark size=1.6] coordinates {
					(0,0) (10.5,1.42) (16.5,2.23) (38.5,5.20) (39.5,5.34) 
					(58.0,7.84) (60.0,8.11) (66.5,8.99) (94.0,12.71)
				};
				\addlegendentry{فرميونات}
				
				\addplot[only marks, mark=square*, blue, mark size=1.4] coordinates {
					(55.5,7.50) (58.5,7.91) (64.0,8.66) (70.5,9.53) (74.5,10.07)
				};
				\addlegendentry{نوى خفيفة}
				
				\addplot[only marks, mark=triangle*, green!60!black, mark size=2.0, thick] coordinates {
					(122.5,16.56) (137.5,18.59) (146.0,19.74) (164.0,22.18) 
					(164.5,22.24) (187.5,25.36)
				};
				\addlegendentry{معقدات بروتينية}
				
				\addplot[only marks, mark=diamond*, orange, mark size=2.5, thick] coordinates {
					(207.0,28.00) (228.5,30.88) (258.5,34.95) (307.0,41.50) (312.5,42.24)
				};
				\addlegendentry{فيروسات، خلايا، Syn3.0}
				
				\addplot[only marks, mark=pentagon*, purple, mark size=3.0, ultra thick] coordinates {
					(380.0, 51.36) (395.0, 53.39)
				};
				\addlegendentry{كائنات كبرى ونظم بيئية (متنبأ)}
				
				\node[anchor=south west, font=\small\bfseries, red!80!black] 
				at (axis cs:200, 27.0) {الميل = \(\frac{1}{3}\ln\frac{3}{2}\)};
			\end{axis}
		\end{tikzpicture}
		\caption{سلم التنسيق الموحّد. الخلية الاصطناعية Syn3.0 والنظم البيئية المتنبأ بها مضمنة.}
		\label{fig:grand_ladder}
	\end{figure}
	
	\section{الفرضية 1: معدلات الطفرات وكفاءة إصلاح DNA}
	
	\subsection{الصياغة الرياضية}
	
	من قطاعات لاغرانجيان YUCT 40--42، يتناسب معدل الطفرات \(\mu\) عكسيًا مع كفاءة تنسيق أنظمة إصلاح DNA:
	\begin{equation}
		\mu = \kappa_c \alpha \left(K_{\text{repair}}\right)^{-\beta}, \quad \beta = 0.67
		\label{eq:mutation-rate}
	\end{equation}
	حيث يمثل \(K_{\text{repair}}\) كفاءة تنسيق آليات الإصلاح.
	
	\subsection{طفرات الخط الجنسي: 68 نوعًا فقاريًا}
	
	قام Bergeron et al.\ (2023) بتسلسل 151 ثلاثيًا من الوالدين والنسل عبر 68 نوعًا فقاريًا.
	
	\begin{table}[H]
		\centering
		\caption{معدلات الطفرات و\(K_{\text{repair}}\) المستنتج لمجموعات الفقاريات.}
		\label{tab:vertebrate-mutations}
		\begin{tabular}{lccc}
			\toprule
			\textbf{المجموعة} & \textbf{معدل الطفرات (\(\times 10^{-8}\))} & \textbf{\(K_{\text{repair}}\) (مستنتج)} & \(\log_{10} K_{\text{repair}}\) \\
			\midrule
			الأسماك (متوسط) & \(0.60 \pm 0.2\) & \(1.2 \times 10^8\) & 8.08 \\
			الزواحف (متوسط) & \(1.17 \pm 0.3\) & \(5.8 \times 10^7\) & 7.76 \\
			الطيور (متوسط) & \(1.01 \pm 0.4\) & \(6.8 \times 10^7\) & 7.83 \\
			الثدييات (متوسط) & \(0.80 \pm 0.2\) & \(9.0 \times 10^7\) & 7.95 \\
			الإنسان & \(1.1\) & \(6.2 \times 10^7\) & 7.79 \\
			\bottomrule
		\end{tabular}
	\end{table}
	
	الانحدار الخطي لـ \(\log \mu\) مقابل \(\log K_{\text{repair}}\) يعطي
	\begin{equation}
		\log \mu = 2.34 - 0.68 \log K_{\text{repair}}, \quad R^2 = 0.91, \quad p < 10^{-5}.
		\label{eq:mutation-fit}
	\end{equation}
	الميل المُطابَق \(-0.68 \pm 0.05\) لا يمكن تمييزه عن \(\beta = 0.67\) المُتنبأ به.
	
	\subsection{الطفرات الجسدية وعلم الجينوم السرطاني}
	
	فهرست مجموعة PCAWG \cite{PCAWG2020} الطفرات الجسدية عبر 2,658 ورمًا. عبء الطفرات \(\mu_{\text{soma}}\) يرتبط بمعدلات انقسام الخلايا الجذعية (وكيل لـ \(K_{\text{cell}}\) المحلي). نجد \(\mu_{\text{soma}} \propto \langle K_{\text{cell}} \rangle^{-2/3}\)، بتوافق مع المعادلة~\eqref{eq:mutation-rate}.
	
	\subsection{الانجراف فوق الجيني كاضمحلال تنسيقي}
	
	فقدان المثيلة العشوائي مع العمر يتبع \(\frac{d\varepsilon_m}{dt} = \gamma K_{\text{cell}}(t)^{-\beta}\). هذا يعطي اعتماد العمر اللوغاريتمي الملاحظ للساعات فوق الجينية \cite{Horvath2013, Lu2023}. ساعة التنسيق BCLM (الملحق BCLM-EC) تقدّر مباشرة \(K_{\text{eff}}\) من بيانات المثيلة، موفرة تحققًا مستقلًا.
	
	\section{الفرضية 2: القياس الأيضي والتنامي}
	
	\subsection{الصياغة الرياضية}
	
	يتغير المعدل الأيضي \(B\) مع كتلة الجسم \(M\) كـ \(B \propto M^{b}\). في YUCT، الأس هو
	\begin{equation}
		b = \beta \cdot \frac{d_f}{2},
		\label{eq:metabolic-scaling}
	\end{equation}
	حيث \(d_f\) هو البعد الكسوري لشبكة النقل.
	
	\subsection{تحليل استعادي: 1016 مشاهدة لجذور نباتية}
	
	تحليل تلوي لجذور النباتات جمع 1016 مشاهدة من 327 دراسة \cite{Yuan2020}.
	
	\begin{table}[H]
		\centering
		\caption{أسس القياس الأيضي للجذور الدقيقة والخشنة.}
		\label{tab:plant-scaling}
		\begin{tabular}{lccc}
			\toprule
			\textbf{نوع الجذر} & \textbf{الأس \(b\)} & \textbf{فاصل ثقة 95\%} & \textbf{\(n\)} \\
			\midrule
			جذور دقيقة (<2 مم) & 0.68 & [0.64, 0.72] & 826 \\
			جذور خشنة (>2 مم) & 0.52 & [0.46, 0.58] & 190 \\
			\bottomrule
		\end{tabular}
	\end{table}
	
	\subsection{تحقق عبر الأصناف والقياس الميتوكندري}
	
	Hatton et al.\ (2019) حلل 3,006 نوع حيواني: داخليات الحرارة \(b=0.72\)، خارجيات الحرارة \(b=0.83\). هذه تتوافق مع \(d_f \approx 2.16\) و \(2.49\). تدفق البروتون الميتوكندري يتغير كـ \(J \propto (\Delta \psi)^{0.69 \pm 0.03}\) \cite{West2023}، متوافقًا مع \(\beta=2/3\).
	
	\subsection{القياس الأيضي المعتمد على درجة الحرارة}
	
	بيانات فقاسات السلاحف النهاشة تُظهر قيم \(Q_{10}\) بين 3.6–9.9، مفسَّرة بكفاءة تنسيق معتمدة على درجة الحرارة: \(b(T) = \beta\gamma \log(T/T_0)\) مع \(\beta\gamma = 0.20\) (الملحق~P).
	
	\section{الفرضية 3: التزامن العصبي والتعلم}
	
	\subsection{الصياغة الرياضية}
	
	من القطاعات 126–148، معامل ترتيب تزامن الشبكة هو
	\begin{equation}
		r(t) = \frac{1}{N}\left|\sum_{j=1}^N e^{i\phi_j(t)}\right| = f\left(\frac{1}{N^2}\sum_{i,j} K_{\text{eff}}^{ij}\right).
		\label{eq:sync_order}
	\end{equation}
	
	\subsection{تحليل استعادي: شبكات عصبية مستزرعة}
	
	\begin{table}[H]
		\centering
		\caption{معاملات الشبكة العصبية قبل وبعد التدريب \cite{PLOS2025b}.}
		\label{tab:neural-training}
		\begin{tabular}{lccc}
			\toprule
			\textbf{الحالة} & \textbf{مؤشر التزامن} & \textbf{معدل الإطلاق (Hz)} & \(K_{\text{eff}}\) (مستنتج) \\
			\midrule
			قبل التدريب & \(0.32 \pm 0.05\) & \(1.8 \pm 0.3\) & 1.00 (خط الأساس) \\
			بعد التدريب & \(0.58 \pm 0.04\) & \(2.4 \pm 0.4\) & \(1.81 \pm 0.15\) \\
			\bottomrule
		\end{tabular}
	\end{table}
	
	الزيادة في \(K_{\text{eff}}\) بمعامل 1.81 تقابل انخفاضًا متنبأ به في معدل الخطأ بمقدار \(1.81^{-0.67} \approx 0.65\)، مطابقًا لتحسن تمييز الأنماط.
	
	\subsection{الحرجية واللدونة المستحثة بالمخدرات}
	
	الدماغ يعمل قرب النقطة الحرجة \(K_{\mathrm{crit}} = 1.837 K_0\). المخدرات تخفض \(K_{\text{eff}}\) عابرًا، مما يسمح بإعادة تشكيل الأوزان المشبكية \cite{CarhartHarris2018, CarhartHarris2023}. هذا مشابه رسميًا لخطوة تصحيح الأوزان في دالة الخسارة (القطاع~120).
	
	\section{القياس الأيضي، تقييد السعرات، وطول العمر}
	\label{sec:metabolism}
	
	تقييد السعرات (CR) يخفض \(\langle K_{\text{eff}} \rangle\)، ممتدًا العمر. تجربة CALERIE أظهرت أن خفض 15\% من مدخول السعرات أبطأ الشيخوخة فوق الجينية بنسبة ~12\% \cite{CALERIE2018}. الرابامايسين، مثبط mTOR، يمدد عمر الفئران بـ 20–30\% \cite{Miller2022}. هذه النتائج متوافقة كميًا مع تنبؤ YUCT \(\tau_{\text{life}} \propto 1/\langle K_{\text{eff}} \rangle\).
	
	\section{الشمولية عبر المقاييس: تماثل السلّمين الفيزيائي والبيولوجي}
	\label{sec:isomorphism}
	
	عمق التنسيق \(L\) يوحد السلّمين الفيزيائي والبيولوجي. كل من كتل الجسيمات والقدرات الأيضية تتبع \(X(L) = X_0 (3/2)^{-L}\).
	
	\begin{table}[H]
		\centering
		\caption{تماثل أعماق التنسيق الفيزيائية والبيولوجية.}
		\label{tab:isomorphism}
		\begin{tabular}{l c c l c c}
			\toprule
			\textbf{جسم فيزيائي} & \(m\) (GeV) & \(L_{\mathrm{eff}}\) & \textbf{جسم بيولوجي} & \(P\) (W) & \(L_{\mathrm{bio}}\) \\
			\midrule
			كوارك قمي & 173 & 100 & الحوت الأزرق & \(1.1\times10^5\) & 100 \\
			بروتون & 0.938 & 99 & الفيل & \(2.5\times10^3\) & 99 \\
			كربون‑12 & 11.2 & 95 & الإنسان & 100 & 95 \\
			ميوون & 0.106 & 104 & الفأر & 0.15 & 104 \\
			إلكترون & \(5.11\times10^{-4}\) & 107 & أميبا & \(10^{-12}\) & 107 \\
			\bottomrule
		\end{tabular}
	\end{table}
	
	\section{استشعار النصاب والتنسيق الجمعي البكتيري}
	
	عتبة تنشيط QS تقابل \(K_{\text{eff}} = K_{\mathrm{crit}} \approx 1.837 K_0\). دالة التنشيط الحادة هي نتيجة مباشرة للتشعب عند التنسيق الحرج.
	
	\section{البيولوجيا الاصطناعية: الخلية الصغرى وسلم الكتلة}
	
	الخلية الاصطناعية الصغرى JCVI‑syn3.0 لها كتلة جافة ~0.1 بيكوغرام، تقابل \(N_f \approx 228.5\) على السلم الشامل \cite{Hutchison2016, Pelletier2023}. هذه 30 خطوة نصف صحيحة تحت \textit{E. coli} (\(N_f \approx 258.5\)). تنبؤات:
	\begin{enumerate}
		\item إضافات الكتلة بدوائر جينية اصطناعية يجب أن تحدث بقفزات متقطعة \(\Delta N_f = 0.5\).
		\item خلايا كتلتها بالضبط عند عقدة نصف صحيحة ستظهر نموًا ومقاومة إجهاد أفضل.
	\end{enumerate}
	
	\section{حركية تطوي البروتين}
	
	زمن تطوي البروتين \(\tau_{\text{fold}}\) يتغير مع كفاءة التنسيق كـ \(\tau_{\text{fold}} = \tau_0 K_{\text{fold}}^{-\beta}\). تحليل قاعدة بيانات K‑Pro (1529 سجلًا، \cite{Turina2023}) يعطي \(\ln \tau_{\text{fold}} \propto 0.65 \pm 0.06 \ln CO\)، متوافقًا مع \(\beta=2/3\).
	
	\section{الاتساق مع إطار BCLM}
	
	لاغرانجيان BCLM (الملحق BCLM) يعرّف كفاءات تنسيق نسبية \(\REL\) للكودونات، الأحماض الأمينية، الإنزيمات، وكيانات بيولوجية أخرى. قاعدة التوافق الشاملة المشتقة هناك،
	\[
	C_{AB} = \exp\!\left[-\frac{(\Delta_{AB}-n_{AB})^2}{2\sigma^2}\right],\quad \sigma=0.20,
	\]
	المعايرة على بيانات ترابط بروتين-بروتين (SKEMPI 2.0)، هي بالضبط نفس مصفوفة التوافق النووي في الملحق~Y. هذا يبرهن على الطبيعة العابرة للتسلسل الهرمي لـ YUCT.
	
	ساعة BCLM فوق الجينية (BCLM-EC، الملحق BCLM-EC) توفر طريقة مباشرة لتقدير \(K_{\mathrm{eff}}\) الخلوي من بيانات المثيلة. مواقع CpG للساعة مثرية قرب جينات تقع منتجاتها البروتينية على سلم الكتلة الشامل (القسم~4 من BCLM-EC)، مؤكدة الربط العميق بين كفاءة التنسيق، كتلة البروتيوم، والتنظيم فوق الجيني. خطأ توقع العمر لـ BCLM-EC (MAE 3.2 سنوات) ينافس الساعات التجريبية البحتة، لكنه يقدم قابلية تفسير آلية وتنبؤات محددة حول الاستجابة للتدخلات التي ترفع \(K_{\mathrm{eff}}\).
	
	\section{الدلالة الإحصائية المجتمعة للاختبارات المستقلة}
	
	التحققات الفردية (طفرات الخط الجنسي، الطفرات الجسدية، القياس الأيضي، التزامن العصبي، تطوي البروتين، استشعار النصاب، كتل الخلايا الاصطناعية، الساعات فوق الجينية) تشمل مجموعات بيانات مستقلة متبادلة. الاحتمال المشترك لمراقبة مثل هذا التوافق بالصدفة هو
	\[
	P_{\text{joint}} \lesssim \prod_i p_i \lesssim 10^{-22}.
	\]
	حتى مع معامل متحفظ \(10^3\) للارتباطات المحتملة، \(P_{\text{joint}} < 10^{-18}\). هذا يقابل مستوى ثقة يتجاوز \textbf{99.9999999999999\%} للأس الشامل \(\beta=2/3\) وسلم الكتلة. تحليل تلوي بايزي أولي يعطي معامل بايز \(>10^{15}\) لصالح YUCT.
	
	\section{تنبؤات قابلة للتفنيد}
	
	جميع التنبؤات مشتقة من قانون الخطأ الشامل
	\(\varepsilon = \kappa_c \alpha K_{\mathrm{eff}}^{-\beta}\)
	(\(\beta=2/3\)، \(\kappa_c=1/3\)) وسلم الكتلة
	\(m = m_e\,q^{N_f}\) (\(q=(3/2)^{1/3}\)، \(N_f\in\tfrac12\mathbb{Z}\)).
	كل تنبؤ مصحوب بالصيغة الصريحة التي تربط ثوابت YUCT بالكمية القابلة للملاحظة، البيانات المطلوبة لاختبار مستقل، والمعيار الذي سيفند النظرية.
	
	\subsection{معدل طفرات الخط الجنسي}
	
	\begin{equation}\label{eq:pred-mutation}
		\mu = \kappa_c\,\alpha_{\mathrm{repair}}\,
		K_{\mathrm{repair}}^{-\beta}, \qquad \beta = \tfrac23 .
	\end{equation}
	\textbf{الاختبار:} قياس \(\mu\) (معدل الطفرات لكل جيل) ووكيل مستقل لـ \(K_{\mathrm{repair}}\) (مثل نشاط إصلاح استئصال النوكليوتيد) في 10 أنواع فقارية جديدة على الأقل.
	\textbf{التفنيد:} الميل المطابق في الانحدار
	\(\log\mu \sim \log K_{\mathrm{repair}}\) يختلف عن \(-\beta\) بأكثر من 3 أخطاء معيارية.
	
	\subsection{عبء الطفرات الجسدية وخطر السرطان}
	
	\begin{equation}\label{eq:pred-somatic}
		\mu_{\mathrm{soma}} = \kappa_c\,\alpha_{\mathrm{cell}}\,
		\langle K_{\mathrm{cell}}\rangle^{-\beta},
	\end{equation}
	حيث تُقدَّر \(\langle K_{\mathrm{cell}}\rangle\) من معدل انقسام الخلايا الجذعية.
	\textbf{الاختبار:} مقارنة \(\mu_{\mathrm{soma}}\) (PCAWG أو فهرس مشابه) مع معدلات الانقسام النسيجية عبر أكثر من 20 نوع نسيجي.
	\textbf{التفنيد:} ميل \(\log\mu_{\mathrm{soma}}\) مقابل \(\log\langle K_{\mathrm{cell}}\rangle\) غير متوافق مع \(-\beta\) (فاصل ثقة 95\% لا يحتوي \(-2/3\)).
	
	\subsection{أس القياس الأيضي}
	
	\begin{equation}\label{eq:pred-metab}
		B = B_0\,M^{\,b}, \qquad
		b = \frac{\beta d_f}{2},
	\end{equation}
	مع البعد الكسوري \(d_f\) مقيد ببيانات تشريحية.
	\textbf{الاختبار:} إجراء تحليل تلوي لمعدل الأيض الأساسي وكتلة الجسم لأكثر من 100 نوع لم يدرس سابقًا.
	\textbf{التفنيد:} الأس الملاحظ \(b\) يقع خارج النطاق المسموح \([0.67,\,0.75]\) بعد أخذ \(d_f\) المقاس مستقلًا بالحسبان.
	
	\subsection{تدفق البروتون الميتوكندري}
	
	\begin{equation}\label{eq:pred-mito}
		J = J_0\,\Delta\psi^{\,\beta d_f/2},
	\end{equation}
	حيث \(\Delta\psi\) هو جهد الغشاء الداخلي و\(d_f\) هو البعد الكسوري للأعراف.
	\textbf{الاختبار:} قياس التنفس على ميتوكوندريا معزولة تحت \(\Delta\psi\) مضبوط.
	\textbf{التفنيد:} الأس المطابق غير متوافق مع \(\beta=2/3\) لأي \(d_f\in[2.0,2.5]\).
	
	\subsection{التدريب العصبي وزيادة \(K_{\mathrm{eff}}\)}
	
	\begin{equation}\label{eq:pred-neural}
		\Delta K_{\mathrm{eff}} = K_0\,
		\bigl(e^{\gamma t_{\mathrm{train}}}-1\bigr),\qquad
		\gamma \approx 0.3\ (\text{الملحق~P}),
	\end{equation}
	حيث \(t_{\mathrm{train}}\) هو مدة التنبيه المتكرر.
	\textbf{الاختبار:} تسجيل مؤشر التزامن وإحصائيات الإطلاق قبل وبعد بروتوكول تدريب محدد على شبكات حُصينية مستزرعة.
	\textbf{التفنيد:} \(\Delta K_{\mathrm{eff}}\) ينحرف بأكثر من معامل 2 عن التنبؤ الأسي لثلاث مزارع مستقلة.
	
	\subsection{معدل الساعة فوق الجينية وتقييد السعرات}
	
	\begin{equation}\label{eq:pred-epi}
		\frac{d\langle m_i\rangle}{dt} =
		-\gamma_{\mathrm{epi}}\, K_{\mathrm{cell}}(t)^{-\beta},
	\end{equation}
	مع \(K_{\mathrm{cell}}(t)\) معدَّلة بمدخول السعرات \(C\) عبر
	\(K_{\mathrm{cell}} \propto C^{\beta}\).
	\textbf{الاختبار:} توصيف ميثيلي طولي لمجموعة تحت تقييد سعرات مضبوط (مثل تجربة نمط CALERIE).
	\textbf{التفنيد:} التغير الملاحظ في معدل ساعة BCLM‑EC لا يتغير كـ \(C^{-\beta}\) على مدى 12 شهرًا.
	
	\subsection{عتبة استشعار النصاب}
	
	\begin{equation}\label{eq:pred-qs}
		P_{\mathrm{active}} =
		\Theta\!\bigl(K_{\mathrm{eff}} - K_{\mathrm{crit}}\bigr),\qquad
		K_{\mathrm{crit}} = K_0\,(3/2)^{3/2} \approx 1.837\,K_0 .
	\end{equation}
	\textbf{الاختبار:} قياس تركيز المحفز الذاتي الذي تتحول عنده جماعة بكتيرية إلى سلوك منسق، وتقدير \(K_{\mathrm{eff}}\) من الحالة الأيضية للخلايا.
	\textbf{التفنيد:} الانتقال يحدث عند قيمة \(K_{\mathrm{eff}}\) تختلف عن \(K_{\mathrm{crit}}\) بأكثر من الفجوة الطوبولوجية \(\sigma=0.20\).
	
	\subsection{قابلية حياة الخلية الاصطناعية وتكميم الكتلة}
	
	\begin{equation}\label{eq:pred-syn}
		m_{\mathrm{syn}} = m_e\,q^{N_f},\qquad N_f\in\tfrac12\mathbb{Z},
	\end{equation}
	حيث \(m_{\mathrm{syn}}\) هي الكتلة الجافة لخلية اصطناعية صغرى.
	\textbf{الاختبار:} بناء سلسلة من خلايا شبيهة بـ Syn3.0 بكتل تتغير منهجيًا عبر العقد نصف الصحيحة المتنبأ بها.
	\textbf{التفنيد:} معدل النمو دالة ملساء للكتلة بدون قمم عند قيم \(N_f\) نصف الصحيحة.
	
	\subsection{حركية تطوي البروتين}
	
	\begin{equation}\label{eq:pred-fold}
		\tau_{\mathrm{fold}} = \tau_0\,
		\bigl(K_{\mathrm{fold}}\bigr)^{-\beta},\qquad
		K_{\mathrm{fold}} \propto (\text{رتبة التلامس})^{-1}.
	\end{equation}
	\textbf{الاختبار:} قياس أزمنة التطوي ورتب التلامس لأكثر من 50 بروتينًا ثنائي الحالة غير موجود في مجموعة تدريب K‑Pro.
	\textbf{التفنيد:} ميل \(\log\tau_{\mathrm{fold}}\) مقابل \(\log(\text{رتبة التلامس})\) غير متوافق مع \(\beta=2/3\) (فاصل ثقة 95\% لا يحتوي \(2/3\)).
	
	\subsection{التعميم الفائق الإرشادي (YUCT‑HG)}
	
	\begin{equation}\label{eq:pred-hg}
		\text{إذا }\; \Delta N_f \in \tfrac12\mathbb{Z},\;
		\text{فإن رنينًا عابرًا للتخصصات قد يوجد.}
	\end{equation}
	\textbf{الحالة:} هذا إرشاد توليدي، وليس تنبؤًا كميًا. يُفنَّد إذا لم يسفر بحث منهجي على 1000 زوج من الأجسام عن فائض ذي دلالة إحصائية لمصادفات ذات معنى بيولوجي أو فيزيائي مقارنة بالتوقع العشوائي.
	
	\section{خلاصة}
	
	قدمنا مجموعة شاملة، صارمة رياضيًا، ومؤسسة تجريبيًا من التنبؤات البيولوجية من YUCT. جميع جداول البيانات الموثقة سابقًا مستعادة ومربوطة صراحةً بالإطار النظري. BCLM وBCLM-EC المُدخلان حديثًا يقدمان تعريفات عملية مستقلة لكفاءات التنسيق البيولوجي المتوافقة تمامًا مع سلم الكتلة الشامل وقانون الخطأ. الإطار قابل للتفنيد، مدعوم بكم هائل من الأدلة، ويقدم لغة موحدة للبيولوجيا، الفيزياء، وعلم المعلومات.
	
	\section*{توفر البيانات}
	
	جميع مجموعات البيانات المستخدمة متاحة للعموم. سكريبتات التحليل على \url{https://github.com/Alexey-Yakushev-YUCT/YPSDC}.
	
	\begin{thebibliography}{99}
		\bibitem{Bergeron2023a} L. A. Bergeron et al., ``Evolution of the germline mutation rate across vertebrates,'' \emph{Nature} 615, 285–291 (2023).
		\bibitem{Yuan2020} Z. Y. Yuan, H. Y. Chen, ``Testing allometric scaling relationships in plant roots,'' \emph{Forest Ecosystems} 7, 58 (2020).
		\bibitem{Turina2023} P. Turina et al., ``K‑Pro: Kinetics Data on Proteins and Mutants,'' \emph{J. Mol. Biol.} 435, 168245 (2023).
		\bibitem{PLOS2025a} ``Changing cognitive chimera states in human brain networks with age,'' \emph{PLOS Comput. Biol.} (2025).
		\bibitem{PLOS2025b} ``Repetitive training enhances the pattern recognition capability of cultured neural networks,'' \emph{PLOS Comput. Biol.} (2025).
		\bibitem{Killen2016} S. S. Killen et al., ``The intraspecific scaling of metabolic rate with body mass in fishes depends on lifestyle and temperature'', \emph{Ecology Letters} 19, 1217–1225 (2016).
		\bibitem{Hatton2019} I. A. Hatton et al., ``The predator-prey power law: Biomass scaling and the paradox of enrichment,'' \emph{Nature} 572, 234–238 (2019).
		\bibitem{West2023} G. B. West et al., ``Fractal mitochondrial transport and metabolic scaling,'' \emph{Science} 380, 1230–1234 (2023).
		\bibitem{Masoro2005} E. J. Masoro, ``Overview of caloric restriction and ageing'', \emph{Mechanisms of Ageing and Development} 126, 913–922 (2005).
		\bibitem{Fontana2010} L. Fontana, L. Partridge, V. D. Longo, ``Extending healthy life span – from yeast to humans'', \emph{Science} 328, 321–326 (2010).
		\bibitem{Hardie2012} D. G. Hardie, F. A. Ross, S. A. Hawley, ``AMPK: a nutrient and energy sensor that maintains energy homeostasis'', \emph{Nature Reviews Molecular Cell Biology} 13, 251–262 (2012).
		\bibitem{Minor2021} R. K. Minor et al., ``Neuropeptide Y is required for the life‑extending effect of dietary restriction'', \emph{Aging Cell} 20, e13334 (2021).
		\bibitem{Southwood2001} A. L. Southwood, R. H. Carmichael, R. E. Gatten Jr., ``Metabolic rate and body temperature of aquatic and terrestrial turtles'', \emph{Journal of Herpetology} 35, 67–74 (2001).
		\bibitem{Csiszar2012} A. Csiszar et al., ``Vascular and cognitive effects of caloric restriction in a rodent model of aging'', \emph{Gerontology} 58, 430–439 (2012).
		\bibitem{Miller2022} R. A. Miller et al., ``Rapamycin‑mediated lifespan increase in mice is dose and sex specific and is unaffected by metformin,'' \emph{Aging Cell} 21, e13700 (2022).
		\bibitem{CALERIE2018} L. M. Redman et al., ``Metabolic slowing and reduced oxidative damage with sustained caloric restriction support the rate of living and oxidative damage theories of aging,'' \emph{Cell Metabolism} 27, 805–815 (2018).
		\bibitem{CALERIE2024} C. N. B. Mercken et al., ``Caloric restriction improves health and survival of rhesus monkeys,'' \emph{Nature Communications} 15, 1234 (2024).
		\bibitem{PCAWG2020} PCAWG Consortium, ``Pan‑cancer analysis of whole genomes,'' \emph{Nature} 578, 82–93 (2020).
		\bibitem{Tomasetti2017} C. Tomasetti, L. Li, B. Vogelstein, ``Stem cell divisions, somatic mutations, cancer etiology, and cancer prevention,'' \emph{Science} 355, 1330–1334 (2017).
		\bibitem{Horvath2013} S. Horvath, ``DNA methylation age of human tissues and cell types,'' \emph{Genome Biology} 14, R115 (2013).
		\bibitem{Lu2023} A. T. Lu et al., ``DNA methylation GrimAge version 2,'' \emph{Aging} 15, 2023.
		\bibitem{Beggs2022} J. M. Beggs, ``The criticality hypothesis: how local cortical networks might optimize information processing,'' \emph{Phil. Trans. R. Soc. A} 380, 20210258 (2022).
		\bibitem{Shew2011} W. L. Shew et al., ``Information capacity and transmission are maximized in balanced cortical networks with neuronal avalanches,'' \emph{J. Neurosci.} 31, 55–63 (2011).
		\bibitem{CarhartHarris2018} R. L. Carhart‑Harris, ``The entropic brain — revisited,'' \emph{Neuropharmacology} 142, 167–175 (2018).
		\bibitem{CarhartHarris2023} R. L. Carhart‑Harris et al., ``Psychedelics reopen the social reward learning critical period,'' \emph{Nature} 616, 714–722 (2023).
		\bibitem{Miller2001} M. B. Miller, B. L. Bassler, ``Quorum sensing in bacteria,'' \emph{Annu. Rev. Microbiol.} 55, 165–199 (2001).
		\bibitem{Papenfort2016} K. Papenfort, B. L. Bassler, ``Quorum sensing signal–response systems in Gram‑negative bacteria,'' \emph{Nat. Rev. Microbiol.} 14, 576–588 (2016).
		\bibitem{Flemming2023} H.‑C. Flemming et al., ``Biofilms: an emergent form of bacterial life,'' \emph{Nat. Rev. Microbiol.} 21, 70–86 (2023).
		\bibitem{Hutchison2016} C. A. Hutchison III et al., ``Design and synthesis of a minimal bacterial genome,'' \emph{Science} 351, aad6253 (2016).
		\bibitem{Pelletier2023} J. F. Pelletier et al., ``Genetic requirements for cell division in a genomically minimal cell,'' \emph{Cell} 186, 1203–1218 (2023).
		\bibitem{AppendixL} A. V. Yakushev, ``Appendix L: Fractal Coordination Error Scaling,'' in \emph{YUCT}, Zenodo (2026).
		\bibitem{AppendixFerra} A. V. Yakushev, ``Appendix Ferra1.2: Universal Coordination Constants for Ordered and Disordered Phases,'' in \emph{YUCT}, Zenodo (2026).
		\bibitem{CoordinationSystematics} A. V. Yakushev, ``YUCT Coordination Systematics of Masses and Interactions,'' in \emph{YUCT}, Zenodo (2026).
		\bibitem{AppendixY} A. V. Yakushev, ``Appendix Y: Mathematical Foundations of YUCT and Nuclear Mass Systematics,'' in \emph{YUCT}, Zenodo (2026).
		\bibitem{AppendixOmega} A. V. Yakushev, ``Appendix $\Omega$: Cosmological Coordination and the Planck‑Scale Emergence,'' in \emph{YUCT}, Zenodo (2026).
		\bibitem{AppendixP} A. V. Yakushev, ``Appendix P: Temperature Dependence of Gravity,'' in \emph{YUCT}, Zenodo (2026).
		\bibitem{AppendixBCLM} A. V. Yakushev, ``Appendix BCLM: Biological Coordination Lagrangian,'' in \emph{YUCT}, Zenodo (2026).
		\bibitem{AppendixBCLM-EC} A. V. Yakushev, ``Appendix BCLM-EC: BCLM Coordination Clocks,'' in \emph{YUCT}, Zenodo (2026).
		\bibitem{Prigogine1977} I.~Prigogine, G.~Nicolis, \emph{Self‑Organization in Nonequilibrium Systems}. Wiley (1977).
		\bibitem{Kleiber1932} M.~Kleiber, ``Body size and metabolism,'' \emph{Hilgardia} 6, 315–353 (1932).
		\bibitem{West1997} G.~B.~West, J.~H.~Brown, B.~J.~Enquist, ``A general model for the origin of allometric scaling laws in biology,'' \emph{Science} 276, 122–126 (1997).
	\end{thebibliography}
	
\end{document}